\documentclass[sigconf,nonacm]{acmart}
\settopmatter{printacmref=false}
\renewcommand\footnotetextcopyrightpermission[1]{}
\usepackage{amsmath}
\usepackage{eqannotate}
\title{EqAnnotate Integration in an ACM-Style Article}
\author{Integration Test}
\affiliation{\institution{Synthetic Validation Fixture}\country{Nowhere}}
\email{test@example.com}
\begin{document}
\begin{abstract}
This fixture checks article-scale spacing, narrow-column wrapping, numbering, and manual fallback in an ACM-style two-column paper.
\end{abstract}
\maketitle
\section{Introduction}
Equation annotations must coexist with dense conference prose and normal page building. This paragraph provides ordinary text before a numbered display and checks that the package uses the current column width.
\begin{annotatedequation}[numbered]
y=\eqmark[blue]{base}{f_\theta(x)}+\eqmark[orange]{corr}{\lambda g(x)}
\label{eq:acm}
\eqannotate{base}{Base prediction}
\eqannotate{corr}{A longer correction description that should wrap within the current ACM column}
\end{annotatedequation}
Text after Eq.~\eqref{eq:acm} should remain below the lower annotation band. The next example uses manual placement without any explicit spacing command in the article source.
\eqannotatecolortheme{mono}
\eqannotatecalloutstyle{arrow}
\begin{annotatedequation}
E=\eqmark[green]{kinetic}{\frac12mv^2}+\eqmark[purple]{potential}{V(x)}
\eqannotate{kinetic}{Kinetic term}
\eqannotatemanual[xshift=13mm,yshift=-13mm]{potential}{Manual potential term}
\end{annotatedequation}
Normal prose resumes after the manual annotation. This final paragraph verifies that no overlay reaches into surrounding text and that the display remains visually part of the article rather than an isolated gallery object.
\section{Conclusion}
Both automatic and manual annotation paths should preserve the surrounding article layout.
\end{document}
